\chapter{結論}\label{ch:conclusion}

この章では、2 週間の実験の結果を、最初の問い「LLM の見通しを Black--Litterman モデルに入れると、時価総額加重を上回るか」に向けてまとめ直す。

\section{4 つの問いへの答え}

\subsection{問い 1：LLM の数値予想は BLM のビューになるか ― ならない}

元論文の形（銘柄ごとに直近の日次リターンを渡し、翌月の平均日次リターンを数で答えさせる）では、gpt-4o-mini の予想は入力の値動きをなぞった（相関 0.89）。入力の要約で説明できない部分の IC は $-0.004$ で、95\% 区間の上限は 0.027 だった（H02）。入力を変えても、質問の形を変えても、会話にしても、答えの形を 5 段階にしても、なぞりは消えなかった（実験 A〜G）。文章だけを渡すとなぞりは消えたが、当たる力は確かめられなかった（H05）。

これは、LLM が「直近の動きが続く」と外挿する傾向を持つという先行研究の報告と一致する。数字の処理を LLM に任せると、LLM は数字の中で最も目立つ信号（直近のリターン）を写すだけになる。

BLM の側にも問題があった。回答の分散を $\Omega$ にすると、$\Omega$ が小さすぎて事後の期待リターンの 99.9\% がビューの側に寄り、市場の均衡がほとんど使われなかった。元の実装の最適化（SLSQP）は途中で止まっていて、その位置は入力のわずかな違いで大きく変わった（H01）。元論文の結果は、この計算の作りに強く左右されている。

\subsection{問い 2：LLM は文章から株価の材料を読めるか ― 読める}

決算発表文の gpt-4o-mini の点数は、発表をはさむ 2 取引日の超過リターンを、1 標準偏差あたり +120bp 説明した（実験 2、$t$ 9.4）。この効果は、次のどれを統制しても残った。

\begin{longtable}{P{6.0cm}P{3.0cm}P{5.4cm}}
\toprule
\textbf{統制に加えたもの} & \textbf{LLM の係数} & \textbf{実験} \\
\midrule
\endhead
\bottomrule
\endfoot
数値の驚き（SUE、売上）、辞書、FinBERT、TF-IDF＋リッジ、発表前のリターン、規模、固定効果 & +120bp & 実験 2 \\
＋アナリスト予想との差 & +116bp & 実験 3 \\
＋反応で追加学習した FinBERT & +102bp & 実験 4 \\
＋ChronoGPT の埋め込み＋リッジ & +112bp & 実験 5 \\
＋ChronoGPT の LoRA & +117bp & 実験 10 \\
\end{longtable}

さらに、先読みのない ChronoGPT（版 2015）の埋め込み＋リッジも、2016〜2023 年（+34bp、$t$ 5.5、実験 5）と、2020〜2026 年（+34bp、$t$ 4.8、実験 11）で、従来の NLP を超えて反応を説明した。gpt-4o-mini の知識が及ばない期間でも、10 年前の言葉しか知らないモデルでも、LLM の読み取りは効く。gpt-4o-mini の結果が「答えを覚えていた」だけではないことの、強い傍証である。

\subsection{問い 3：読めた情報で運用成績は上がるか ― 上がらない}

LLM が読めた情報は、発表日の終値までに株価に入っていた。

\begin{itemize}
\item 発表の翌日（day 1）：+3.6bp（$t$ 0.93、実験 8）。95\% 区間の上限でも +11bp で、2 日の反応（約 +100bp）のほとんどは day 0 に出る。
\item day 2〜21 のドリフト：+10bp（$t$ 0.83、実験 2）。アナリスト予想を統制しても +9bp（実験 3）。
\item 月次の BLM：費用の後 +0.10\%／月（31 か月、$t$ 1.22、実験 7）、+0.003\%／月（約 10 年、$t$ 0.04、実験 9）。
\item 日次の BLM：費用の後 年 $-20\%$（回転率 1 日 42\%、実験 8）。
\end{itemize}

ただし、点数を銘柄の間で混ぜた BLM と比べると、本物の点数の BLM は費用の前で月 +0.20\%（$t$ 2.30、実験 7）、+0.15\%（$t$ 1.78、実験 9）上回った。比率を動かすこと自体の損（費用と TE）を除くと、点数の中身に弱い価値が残っている可能性はある。しかしその価値は、比率を動かす費用を払うと消える大きさだった。

\subsection{問い 4：どこを直せば良くなるか ― 単純な直しでは良くならない}

文章の掃除（修正 1）、読む範囲（修正 2）、セクターへのまとめ方（修正 3）、埋め込みの取り方（修正 4）、圧縮（修正 5）を、開発の期間の機械的な規則で試したが、主な切り方 B では残った修正はなかった（実験 11）。今のやり方（直近 92 日の最新の発表を、時価総額加重でセクターにまとめる）が最もよかった。

探索では、ChronoGPT の最後の層ではなく中間の層（52 層のうち 26 層目）の埋め込みが、$R^2$ の増分を 0.49 から 0.80 に上げた（$t$ 2.1〜2.9）。事前登録した比較ではないので、次に確かめる候補にとどまる。

\section{モデルの強さの順位}

同じ標本・同じ統制で比べた、文章のモデルの説明力（統制に加えたときの決定係数の増分、ポイント）は、次の順だった。値の幅は、実験 4・5・10（標本 A、2023-11 以降）、実験 5（標本 B、2016〜2023）、実験 11（切り方 B の最終の期間、2020 以降）の値の範囲である。どの実験の中でも、この順位は入れ替わらなかった。

\begin{longtable}{P{0.6cm}P{5.6cm}P{2.0cm}P{6.2cm}}
\toprule
 & \textbf{モデル} & \textbf{増分} & \textbf{学習} \\
\midrule
\endhead
\bottomrule
\endfoot
1 & gpt-4o-mini（学習なし、ゼロショット） & +2.01〜2.49 & なし。最も大きく、最も新しい \\
2 & FinBERT（全体を反応で追加学習） & +0.84〜1.01 & 全体を反応で \\
3 & ChronoGPT の埋め込み＋リッジ & +0.28〜0.59 & 線形の頭だけ \\
4 & TF-IDF＋リッジ & +0.06〜0.37 & 語の重みを反応で \\
5 & FinBERT（学習なし） & +0.24 & なし \\
6 & ChronoGPT の LoRA & +0.13 & 中まで反応で \\
7 & 辞書 & +0.02 & なし \\
\end{longtable}

\begin{keybox}[順位から読めること]
\begin{enumerate}
\item \textbf{大きく新しい LLM は、学習なしで最も強い}。反応という弱く雑音の多い目的変数に合わせて学習するより、大量の文章で得た一般的な読解力のほうが効く。
\item \textbf{反応で学習させると伸びる}。FinBERT は学習なしの +0.2 から、追加学習で +0.8〜1.0 になった。
\item \textbf{中まで学習させるのは難しい}。ChronoGPT の LoRA は、線形の頭だけのリッジより弱かった。学習の損失はほとんど下がらず、データ（約 1.4 万本）と弱い目的変数に対して、学習する数（約 256 万）が多すぎたと考えられる。
\item \textbf{先読みのない小さな LLM でも、従来の NLP 3 つを合わせたもの（+0.51）を超える}（+0.59、実験 11）。
\end{enumerate}
\end{keybox}

\section{LLM-BLM への示唆}

BLM のビューは、「市場の均衡（時価総額の比率が示す期待リターン）と違う見通し」である。ビューが価値を持つには、その見通しが\textbf{判断の時点で、まだ価格に入っていない}必要がある。この研究で LLM が読めた情報は、価格に入るのが速すぎた。

\begin{figure}[htbp]
\centering
\begin{tikzpicture}[x=12mm,y=8mm]
\draw[->,color=ink2] (0,0) -- (10.5,0) node[right,font=\scriptsize] {時間};
\draw[->,color=ink2] (0,0) -- (0,3.6) node[above,font=\scriptsize] {株価に入った割合};
\draw[line width=1.2pt,color=c1] (0,0.1) -- (1,0.1) -- (1.3,3.0) -- (2.3,3.15) -- (10,3.2);
\draw[dashed,color=c2] (1,-0.1) -- (1,3.4);
\node[font=\scriptsize,color=c2,anchor=south] at (1,3.4) {発表};
\draw[dashed,color=c3!80!black] (2.3,-0.1) -- (2.3,3.4);
\node[font=\scriptsize,color=c3!80!black,anchor=south] at (2.6,3.4) {day 1};
\draw[dashed,color=c7] (6,-0.1) -- (6,3.4);
\node[font=\scriptsize,color=c7,anchor=south] at (6,3.4) {月次の判断日};
\node[font=\scriptsize,align=left,anchor=west] at (2.5,1.6) {ここで LLM が読んでも、\\もう価格に入っている};
\end{tikzpicture}
\caption{LLM の読み取りが価値を持てない理由。情報は発表の当日に価格に入る。月次の判断日に LLM が読んでも、遅すぎる。}
\label{fig:concl-timing}
\end{figure}

LLM をポートフォリオに生かす道は、次のどれかだと考えられる。

\begin{enumerate}
\item \textbf{速さで勝つ}：発表の直後の数分〜数時間に読む。日中の株価が要り、この研究のデータ（日次の終値）では確かめられない。
\item \textbf{市場が読み違えた発表を探す}：平均では価格に入っていても、文章の向きと当日の株価の向きが食い違った発表には、情報が残っているかもしれない（\chapref{ch:open}）。
\item \textbf{期待リターンではなくリスクに使う}：リターンの情報は 1 日で消えるが、ボラティリティは何週間も続く。文章からリスクの変化を読み、共分散や TE の管理に使う。
\item \textbf{時間をかけて効く情報を読む}：10-K のリスクの記述のように、価格への反映がゆっくりな文章を読む（実験 6 の案）。
\end{enumerate}

\section{元論文について}

元論文は、LLM の数値予想を BLM のビューにして、等金額や平均分散を上回ったと報告している。この研究は、gpt-4o-mini・\SP{}・30 か月で同じ形を再現し、その成績が LLM の予測力ではなく、(1) 予想のなぞり（直近の勝ち組を選ぶ）、(2) $\Omega$ が小さすぎて均衡がほぼ使われないこと、(3) 最適化の打ち切り、の組み合わせで説明できることを示した。元論文とはモデル（Llama・Gemma・Qwen）と期間（2024 年後半の約半年）が違うので、元論文の結論が誤りだとは言えない。ただし、その成績を「LLM の見通しの価値」と解釈するには、この研究の H02・H04 のような対照（入力だけのモデル、混ぜた予想）が必要である。

\section{一言で}

\begin{keybox}[結論]
LLM は決算発表文を読める。gpt-4o-mini の読み取りは、数値の驚き、アナリスト予想、従来の NLP、反応で学習したモデルを超えて、発表への株価の反応を説明する。先読みのない LLM でも同じことが言える。しかし、LLM が読めることは、市場も発表の当日に読んでいる。月次や日次の BLM のビューに入れても、費用の後で時価総額加重を上回る成績にはならなかった。
\end{keybox}
